\chapter{翻牌 / Flop}

\section{本章导入}

本章将进入翻牌圈。翻牌不是简单地“中了就打，没中就弃”，而是需要结合牌面结构、双方范围、位置关系和下注目的进行判断。很多新手在翻牌圈犯错，是因为只看自己的手牌强弱，而忽略了整体局面。

翻牌学习的重点，是理解范围优势、坚果优势、持续下注（C-bet）以及不同牌面上的价值下注和诈唬下注逻辑。

\section{核心概念}

本章将介绍牌面结构、干燥牌面、湿润牌面、范围优势（range advantage）、坚果优势（nut advantage）、价值下注、诈唬下注和 C-bet。读者需要学会从“我有什么牌”转向“我的范围在这个牌面上表现如何”。

\section{新手常见误区}

新手常见误区包括任何未成牌都自动放弃、任何顶对都过度保护、所有牌面都机械 C-bet，以及用没有后续计划的诈唬下注制造困难局面。

\section{实战思考框架}

翻牌圈可以先评估牌面是否击中自己的翻前范围，再判断手牌属于价值、摊牌价值、听牌还是空气牌。行动前应明确下注、过牌或跟注的目的。

\section{牌例拆解}

\subsection{牌例一：待补充}

本牌例将用于分析干燥高牌面上的持续下注。

\subsection{牌例二：待补充}

本牌例将用于分析湿润连接牌面上的控池与保护。

\section{本章总结}

翻牌圈是范围思维的第一道关口。读者应避免只凭命中程度行动，而要把牌面、范围和下注目的结合起来。

\section{课后训练}

请选择三个不同翻牌面，分别判断它们更有利于翻前进攻方还是防守方，并写出你的理由。
